% Generated draft: requires mathematical and source review.
\documentclass[11pt]{article}
\usepackage[margin=25mm]{geometry}
\usepackage{fontspec}
\setmainfont{Latin Modern Roman}
\usepackage{amsmath,amssymb,mathtools,unicode-math}
\setmathfont{Latin Modern Math}
\usepackage{enumitem,needspace,xurl,hyperref}
\setlength{\parindent}{0pt}
\setlength{\parskip}{5pt}
\setlength{\emergencystretch}{3em}
\newcommand{\N}{\mathbb{N}}
\newcommand{\Z}{\mathbb{Z}}
\newcommand{\Q}{\mathbb{Q}}
\newcommand{\R}{\mathbb{R}}
\newcommand{\C}{\mathbb{C}}
\newcommand{\F}{\mathbb{F}}
\newcommand{\A}{\mathbb{A}}
\newcommand{\PP}{\mathbb{P}}
\newcommand{\E}{\mathbb{E}}
\newcommand{\eps}{\varepsilon}
\newcommand{\dd}{\,\mathrm{d}}
\begin{document}
\section*{UnsolvedMath ID: 30001024}
\textbf{Source problem code:} catalogue entry
\subsection{Definitions and mathematical statement}
This source poses a morphogenesis research question rather than a specified mathematical conjecture. Phyllotaxis concerns the arrangement of plant organs such as leaves; the source also includes plant ridges and recurring patterns on mollusk shells. Its two linked requests are:
\begin{enumerate}
\item Identify physical or chemical mechanisms that generate the recurring periodic patterns in these plant and mollusk systems.
\item Explain why a relatively small collection of patterns dominates in both groups, often resembling patterns elsewhere in nature and in laboratory experiments.
\end{enumerate}
The word ``periodic'' is retained from the source without imposing an exact spatial lattice, temporal period, or symmetry group on all the biological examples. Likewise, the title does not impose the existence of one mechanism shared by every organism. The question permits comparison of mechanisms and pattern-selection explanations across systems.

No organism class, tissue geometry, growth law, mechanical constitutive relation, chemical kinetics, boundary conditions, measured pattern descriptor, or quantitative criterion of dominance is provided. A precise model-specific formulation would need to specify these data and what counts as explaining the observed recurrence and selection. No reaction--diffusion equation, elastic instability, optimization principle, or universal pattern theorem is selected without source evidence.
\subsection{Short English statement}
What physical or chemical mechanisms create recurring plant and mollusk-shell patterns, and why do only a few patterns predominate, with similarities across nature and laboratory systems? The source gives a broad comparative morphogenesis agenda.
\subsection{Sources}
\begin{itemize}
\item Local UnsolvedMath registry, permanent ID \(30001024\), source code catalogue entry, ``Universal Mechanisms of Biological Pattern Formation.'' No individual proposer is supplied.
\item \emph{Mini-Workshop: The Mathematics of Growth and Remodelling of Soft Biological Tissues} (2009), Oberwolfach Reports, DOI 10.4171/owr/2008/39, pp. 2219--2260: \url{https://doi.org/10.4171/owr/2008/39}.
\item Supplied reference, ``Moulton, Goriely and Chirat, hierarchical mechanical patterns'': \url{https://doi.org/10.1098/rsif.2024.0918}.
\item Supplied literature triage, checked 2026-08-21: status ``open.'' It reports mechanistic explanations for important plant and mollusk pattern classes but calls the broad cross-organism predominance question open. This is supplied dated curation, not independently verified; time-sensitive status and cited claims require independent verification.
\end{itemize}
\end{document}
